\section{Chapitre~\ref{sec:sensors}. Les entrées de la machine}

La structure des capteurs est mesurée directement, par des enregistrements du courant de cellules isolées, des vibrations de la cochlée et par le comptage des cellules de la rétine; la limite de sensibilité et la formule du bruit de grenaille découlent de la physique; que les codes des capteurs soient ajustés pour porter le plus d'information est une hypothèse solidement appuyée.

{\footnotesize
\setlength{\tabcolsep}{3pt}\renewcommand{\arraystretch}{1.15}
\begin{longtable}{>{\raggedright}p{0.5cm}>{\raggedright}p{4.0cm}>{\raggedright}p{5.6cm}>{\raggedright}p{2.3cm}>{\raggedright\arraybackslash}p{1.7cm}}
\toprule
N\textsuperscript{o} & Affirmation & Expérience et ce qui est mesuré & Source & Statut \\
\midrule\endfirsthead
\toprule
N\textsuperscript{o} & Affirmation & Expérience et ce qui est mesuré & Source & Statut \\
\midrule\endhead
1 & Le capteur est un transducteur: le récepteur produit un courant, et la sortie des cellules sensorielles de l'œil et de l'oreille est du glutamate vers le bus \nt{GLUT}; les premières cellules ne produisent pas d'impulsions (fig.~\ref{fig:transducer}) & Les bâtonnets et les cônes de la tortue libèrent un acide aminé excitateur, détecté par des canaux glutamatergiques sur un fragment de membrane excisé. Cellules ciliées internes de la cochlée du rat: les courants dans les terminaisons des fibres nerveuses passent par des récepteurs AMPA du glutamate, et la fréquence des libérations croît avec la dépolarisation progressive de la cellule jusqu'à $150$\,Hz & \cite{CopenhagenJahr1989,GlowatzkiFuchs2002} & mesuré \\
2 & Dans l'obscurité, le capteur de lumière répond à un seul photon par un courant d'environ un picoampère & Électrode à succion sur le segment externe d'un bâtonnet de crapaud: les réponses aux éclairs faibles sont quantifiées, événement de $\approx1$\,pA ($5\%$ du courant d'obscurité) avec une dispersion de $0.2$\,pA, efficacité $0.5$. L'homme signale l'arrivée d'un seul photon plus souvent que le hasard & \cite{Baylor1979,Tinsley2016} & mesuré \\
3 & Le capteur de son détecte un déplacement inférieur au nanomètre & Méthode Mössbauer, membrane basilaire de la cochlée du cobaye au site $16$--$19$\,kHz: au seuil de réponse du nerf auditif, la vitesse de la membrane est d'environ $0.04$\,mm/s, soit un déplacement de $\approx0.35$\,nm (notre conversion). On a mesuré la membrane, et non le faisceau du capteur & \cite{Sellick1982,Hudspeth1989} & mesuré \\
4 & Dans l'obscurité, la précision est limitée par le bruit de grenaille des photons~\eqref{eq:photonnoise}; en lumière faible, l'accumulation est plus longue & La formule découle de la statistique de Poisson. Dans le bâtonnet, le bruit en lumière faible coïncide avec la somme d'événements quantiques indépendants; sur un fond lumineux, la réponse à un éclair est plus petite et plus brève. Le chiffre \og dixièmes de seconde\fg{} n'est pas vérifié séparément ici & \cite{Baylor1979} & théorie \\
5 & La plage d'entrée couvre dix ordres de grandeur pour la lumière et douze pour le son, la fréquence des impulsions moins de trois & Pour le son: la réponse de la membrane basilaire à la fréquence caractéristique croît avec une pente descendant jusqu'à $0.2$\,dB/dB dans la bande $40$--$80$\,dB; la compression reste visible au-delà de $100$\,dB. Les nombres d'ordres de grandeur eux-mêmes sont des estimations classiques, sans source dans le chapitre & \cite{Ruggero1997} & mesuré \\
6 & Réponse du capteur~\eqref{eq:nakarushton}; $\sigma$ suit la luminosité moyenne et fait glisser la courbe le long de l'axe logarithmique (fig.~\ref{fig:adaptation}) & La formule a d'abord été ajustée aux potentiels S de la rétine de poisson. Cônes de la tortue: les réponses obéissent à $V=I V_m/(I+\sigma)$ et couvrent $\sim3.5$ ordres de grandeur de luminosité; après $2$--$3$ minutes de fond, la courbe glisse horizontalement en gardant sa largeur. Lien avec la division par l'activité totale: revue & \cite{NakaRushton1966,NormannPerlman1979,Carandini2012} & mesuré \\
7 & Les capteurs transmettent les changements, et leurs codes sont ajustés pour porter le plus d'information par impulsion (proposition~\ref{prop:sensors}) & Le principe a été proposé en théorie. L'indice le plus fort: chez la mouche, la caractéristique \og contraste $\to$ réponse\fg{} des neurones de la première couche coïncide avec la distribution cumulée des contrastes des scènes naturelles, ce qui maximise l'information & \cite{Barlow1961,Laughlin1981} & hypothèse \\
8 & Les capteurs ON et OFF forment un code à deux fils; la caméra événementielle le reproduit dans le silicium & Revue d'expériences: les voies ON et OFF de la rétine se séparent dès la première synapse et peuvent être supprimées séparément. Caméra: $128\times128$ pixels sans images, plage de $120$\,dB, latence de $15$\,\textmu s & \cite{Schiller1992,Lichtsteiner2008,Gallego2022} & mesuré \\
9 & L'œil compte $\sim10^8$ capteurs de lumière et $\sim10^6$ neurones de sortie: compression d'un facteur $\sim100$ & Comptage sur des rétines humaines entières: en moyenne $92$ millions de bâtonnets et $4.6$ millions de cônes; de $0.7$ à $1.5$ million de cellules de sortie (ganglionnaires) selon les yeux & \cite{Curcio1990,CurcioAllen1990} & mesuré \\
10 & Chez la souris, l'œil a plus de trente canaux de sortie & Souris: imagerie calcique biphotonique de plus de $11\,000$ cellules de sortie et regroupement des réponses, nettement plus de $30$ types fonctionnels. Chez l'homme, le nombre n'est pas mesuré & \cite{Baden2016} & mesuré \\
11 & L'œil du cobaye transmet $\sim875\,000$ bit/s; rapporté à l'œil humain, cela donne $\sim10^7$ bit/s & Cobaye, matrice multiélectrode, mouvement dans des scènes naturelles: $6$--$13$ bit/s par cellule, $\sim875\,000$ bit/s pour $\sim10^5$ cellules de sortie. Pour l'homme ($\sim10^6$ cellules), $\sim10^7$ bit/s: extrapolation des auteurs & \cite{Koch2006} & mesuré \\
12 & L'oreille est un banc de filtres passe-bande avec une carte des fréquences presque logarithmique (fig.~\ref{fig:filterbank}) & Le lieu de vibration maximale de la membrane basilaire dépend de la fréquence; la fonction \og fréquence--lieu\fg{} est presque exponentielle et décrit sous une seule forme les cochlées de l'homme et d'animaux vivants & \cite{Greenwood1990,Robles2001} & mesuré \\
13 & Les résonateurs sont actifs: une partie des capteurs amplifie les vibrations faibles d'un facteur de plusieurs milliers en amplitude ($66$--$76$\,dB); le gain diminue quand le volume augmente & Vibrométrie laser à la base de la cochlée du chinchilla: pour des sons faibles à la fréquence caractéristique, gain de $66$--$76$\,dB par rapport à l'étrier; la mort réduit la sensibilité de $60$--$81$\,dB ($10^3$--$10^4$ fois en amplitude) et supprime la compression. La source de puissance est constituée par les cellules ciliées externes & \cite{Ruggero1997,Robles2001} & mesuré \\
14 & L'amplificateur de l'oreille est placé au bord de l'auto-oscillation & Calcul: près d'une bifurcation de Hopf, compression de la plage, accord fin et sons de combinaison sont inévitables, soit toutes les non-linéarités connues de l'audition. Que la cochlée soit réglée exactement ainsi n'est pas démontré directement & \cite{Eguiluz2000} & hypothèse \\
15 & Aux basses fréquences, les impulsions sont en phase avec le son (chez le cobaye, le verrouillage faiblit à partir de $\sim600$\,Hz et reste perceptible jusqu'à $\sim3.5$\,kHz), et la direction s'en déduit; aux hautes fréquences, il reste le code de position & Nerf auditif du cobaye: le verrouillage de phase commence à décliner vers $600$\,Hz et disparaît au-dessus de $3.5$\,kHz; chez le chat et le singe, la limite est environ une octave plus haut. Différence de temps d'arrivée aux deux oreilles: revue & \cite{PalmerRussell1986,Grothe2010} & mesuré \\
16 & Autres capteurs (tabl.~\ref{tab:sensors}): l'odorat a $\sim400$ types de récepteurs, un par neurone, et un code combinatoire; le goût passe en partie par l'ATP et la sérotonine; le toucher a des capteurs à adaptation rapide et lente; le chaud et le froid sont séparés & Souris: un récepteur reconnaît de nombreuses odeurs, une odeur active de nombreux récepteurs, des odeurs différentes activent des ensembles différents. Sans récepteurs de l'ATP, le nerf gustatif ne répond pas; les bourgeons du goût libèrent de la sérotonine. Enregistrement de fibres cutanées isolées de la main humaine: quatre types selon la vitesse d'adaptation. Le canal du froid est distinct de ceux du chaud & \cite{Mombaerts2004,Malnic1999,Finger2005,Huang2005,JohanssonVallbo1979,McKemy2002} & mesuré \\
\bottomrule
\end{longtable}}
